% =========================================================================
% Capítulo 1 — MODELO DE NEGÓCIOS
% =========================================================================
% Estrutura esperada (espelho do template UNIVALI):
%   1.1 Introdução ao domínio de aplicação
%   1.2 Público alvo
%   1.3 Solução Proposta (com Canvas e SWOT como figuras em images/)
%   1.4 Delimitação do Escopo
%   1.5 Levantamento de produtos existentes (concorrentes)
%       1.5.1 Sympla
%       1.5.2 Eventbrite
%       1.5.3 Facebook Events
%       1.5.4 Instagram (descoberta orgânica)
%       1.5.5 Apps municipais ("Visite Bombinhas", etc.)
%   1.6 Objetivos
%       1.6.1 Objetivos Gerais
%       1.6.2 Objetivos Específicos
%
% Base de informação:
%   - Seção 1, 2, 3, 4, 5 do docs/CONTEXTO_TCC_MONOGRAFIA.md
%   - Imagens necessárias em images/:
%       - ModeloNegocio.pdf  (Business Model Canvas)
%       - AnaliseSWOT.pdf
% =========================================================================

\section{Introdução ao domínio de aplicação}
% TODO: solicitar ao Claude com base na seção 2 do contexto.

\section{Público alvo}
% TODO: solicitar com base na seção 4 (Personas) do contexto.

\section{Solução Proposta}
% TODO: incluir aqui as figuras do Canvas e SWOT (precisam ser desenhadas em
% draw.io / Miro / Canva e exportadas para PDF em images/).

\section{Delimitação do Escopo}
% TODO: delimitar o que o MVP entrega vs. o que fica para o TCC III.
% Base: seção 15 (matriz de status) do contexto.

\section{Levantamento de produtos existentes}
% TODO: subseções para cada concorrente da seção 5 do contexto:
% Sympla, Eventbrite, Facebook Events, Instagram, apps municipais.

\section{Objetivos}
\subsection{Objetivos Gerais}
% TODO

\subsection{Objetivos Específicos}
% TODO: itemize com 5-7 objetivos específicos.
